\documentclass[12pt]{article}

\newif\ifsol
\soltrue
\usepackage{amsthm}

\theoremstyle{plain}
\newtheorem{lemma}{Lemma} 

\theoremstyle{definition}
\newtheorem*{lem_explanation}{Lemma Proof}


% Unicode compatibility
\usepackage{iftex}
\ifPDFTeX
  \usepackage[utf8]{inputenc}
  \usepackage[noTeX]{mmap}
  \usepackage[T1]{fontenc}
\fi
% XeTeX does not support mmap
\ifLuaTeX
  \usepackage{luatex85}
  \usepackage[noTeX]{mmap}
\fi

\usepackage{framed}
%\usepackage{mdframed}
\usepackage[most]{tcolorbox}
\tcbuselibrary{breakable}

\usepackage[most]{tcolorbox}

%\newtcolorbox{mybox}{
%    enhanced,              % Required for advanced break control
%    breakable,             % Allows spanning pages
%    colback=white,
%    colframe=black,
%    sharp corners,
%    boxrule=1pt,
%    % The magic happens here:
%    breakatskip=0pt,
%    bottomrule at break=0pt, % Removes bottom line on the first part
%    toprule at break=0pt,    % Removes top line on the second part
%}


\newtcolorbox{mybox}{
    parbox = true,          % CHANGED BCOZ THIS WAS BREAKING
    enhanced,               % Required for the "at break" keys to work
    breakable,              % Allows the box to span pages
    colback=white,          % Background color
    colframe=black,         % Border color
    sharp corners,          % Square edges like fbox
    boxrule=0.5pt,          % Standard thickness
    % This removes the lines at the break point:
    bottomrule at break=0pt,
    toprule at break=0pt,
    % Optional: Adjust padding to match fbox style
    left=5pt, right=5pt, top=5pt, bottom=5pt
}

\setlength{\textheight}{9in}
\setlength{\textwidth}{7.1in}
\setlength{\evensidemargin}{-0.2in}
\setlength{\oddsidemargin}{-0.2in}
\setlength{\headsep}{30pt}
\setlength{\topmargin}{-0.3in}
\usepackage{amsthm,amsfonts,amsmath,amssymb,caption,xspace,tikz,varwidth}
\usepackage{cancel,subcaption}

\usepackage{hyperref,verbatim}
\usepackage[capitalise, noabbrev]{cleveref}
\newtheorem{theorem}{Theorem}
\theoremstyle{definition}
\newtheorem{definition}{Definition}
\theoremstyle{remark}
\newtheorem{remark}{Remark}

\definecolor{shadecolor}{RGB}{242,242,242}


% \newcounter{defnum}
% \setcounter{defnum}{0}
% \renewenvironment{definition}
%     {\par\addbigskip
%    %\color{red}%
%     \begin{shaded*}
%     \noindent
%     \textbf{Definition.\ }\ignorespaces}
%     {\end{shaded*}
%     \addbigskip}

% \newenvironment{shadeddef}
%     {\begin{shaded*}
%     \noindent
%         \begin{definition}
%         \end{definition}
%     {\end{shaded*}}

\crefname{definition}{Definition}{Definitions}
\crefname{remark}{Remark}{Remarks}



\newcommand{\addmedskip}{\addvspace{\medskipamount}}

\newcommand{\addbigskip}{\addvspace{\bigskipamount}}

\pagestyle{myheadings}
\markboth{BU CAS CS 538.}{BU CAS CS 538.} 
\newcounter{problemnum}
\setcounter{problemnum}{0}
\newenvironment{problem}
     {\addbigskip \setcounter{partnum}{0}
      \noindent\stepcounter{problemnum}\textbf{Problem
                                             \arabic{problemnum}.\ }}
     {\par\addbigskip}


\ifsol
\newenvironment{solution}
     % {\addbigskip
     %  \noindent\textbf{Solution.\ }}
     % {\par\addbigskip}
      {\par\addbigskip
   %\color{red}%
\begin{mybox}\noindent
    \textbf{Solution.\ }\ignorespaces}
    {\end{mybox}
    \addbigskip}
\else
\newenvironment{solution}
     {\comment}
     {\endcomment}
\fi

\newcounter{partnum}
\setcounter{partnum}{0}
\newenvironment{ppart}
     {\addmedskip
      \noindent\stepcounter{partnum}\textbf{(\alph{partnum})}\ }
     {\par\addbigskip}

% \newcommand{\Enc}{\mathrm{Enc}}
% \newcommand{\Dec}{\mathrm{Dec}}
\newcommand{\eqdef}{\buildrel{\scriptstyle{\mathit{def}}}\over{=}}

\newcommand{\mbmod}{\,\%\,}

% \newcommand{\indist}{  \mathrel{\vcenter{\offinterlineskip
%   \hbox{$\sim$}\vskip-.35ex\hbox{$\sim$}\vskip-.35ex\hbox{$\sim$}}}}


\usepackage{enumerate}
\usepackage{enumitem}
\usepackage{mleftright} % for '\mleft' and '\mright' macros
\newcommand{\uprob}[2]{\Pr_{#1}\mleft[\,#2\,\mright]}

% notes
\usepackage{xcolor}
\newcommand{\julia}[1]{\textcolor{red}{\sf {\bf Julia: }#1}}
\newcommand{\leo}[1]{\textcolor{blue}{\sf {\bf Leo: } #1}}

%%%%%%%%%%%%% mike rosulek tex %%%%%%%%%%%%
% %% font stuff

\urlstyle{sf}
\usepackage[T1]{fontenc}
\usepackage{libertine}
%\usepackage[libertine]{newtxmath}
\usepackage[scaled=0.8]{beramono}

\usepackage{microtype}
\usepackage{cmap}


%% colors, styles

\newlength{\myline}
\setlength{\myline}{0.7pt} % pgf "thick"

\colorlet{bg}{white}
\colorlet{fg}{black}
\colorlet{shaded}{black!10}
\colorlet{shadeborder}{black!30}
\definecolor{hlbg}{HTML}{F5F5A4} 
\colorlet{hlfg}{black}
\colorlet{commentcolor}{black!60!white}
\definecolor{errorcolor}{HTML}{a91616}
\definecolor{linkcolor}{HTML}{4b804c}
\definecolor{bitcolor}{HTML}{a91616}

\pagecolor{bg}
\color{fg}

%% highlighting

\usepackage[auto,outline]{contour}

\contourlength{1.2\myline}
\newcommand{\hl}[1]{%
    \relax\ifmmode%
        {}%
        \contour{hlbg}{\textcolor{hlfg}{${} #1 {}$}}%
        {}%
    \else%
        \contour{hlbg}{\textcolor{hlfg}{#1}}%
    \fi%
}

%% codebox stuff 

\usepackage{pipetex}
\pipetexcommand{perl codebox2tex.pl}

\definecolor{boxbordercolor}{HTML}{000000}
\definecolor{boxbgcolor}{HTML}{FFFFFF}

\newcommand{\codebox}[1]{%
    \begin{varwidth}{\linewidth}%
        \upshape%   no slant in definition/theorem statement!
        \begin{tabbing}%
            ~~~\=\quad\=\quad\=\quad\=\kill
            #1
        \end{tabbing}%
    \end{varwidth}%
}

\newcommand{\fcodebox}[1]{%
    \fboxsep=3pt%
    \fcolorbox{boxbordercolor}{boxbgcolor}{\codebox{#1}}%
}

\newcommand{\titlecodebox}[2]{%
    \fboxsep=0pt%
    \fcolorbox{boxbordercolor}{boxbordercolor!10!boxbgcolor}{%
        \begin{varwidth}{\linewidth}%
            \centering%
            \fboxsep=3pt%
            \colorbox{boxbordercolor!10!boxbgcolor}{#1} \\
            \colorbox{boxbgcolor}{\codebox{#2}}%
        \end{varwidth}%
    }%
}

\newcommand{\hltitlecodebox}[2]{%
    \fboxsep=3pt%
    \colorbox{hlbg}{%
        \titlecodebox{#1}{#2}%
    }%
}

\newcommand{\hlcodebox}[1]{%
    \fboxsep=3pt%
    \colorbox{hlbg}{%
        \fcodebox{#1}%
    }%
}

\newcommand{\procheader}[1]{\underline{#1}}
\newcommand{\mycomment}[1]{\textcolor{commentcolor}{\small\textsl{// #1}}}

%% library stuff and math stuff

\renewcommand{\L}{\mathcal{L}}
\newcommand{\lib}[1]{\mathcal{L}_{\textsf{\textup{#1}}}}
\newcommand{\outputs}{\Rightarrow}
\newcommand{\link}{\diamond}

%\newcommand{\indist}{\approxeq}
\renewcommand{\gets}{\twoheadleftarrow}

\renewcommand{\le}{\leqslant}
\renewcommand{\leq}{\leqslant}
\renewcommand{\ge}{\geqslant}
\renewcommand{\geq}{\geqslant}

\newcommand{\pct}{\mathbin{\%}}
\renewcommand{\le}{\leqslant}
\renewcommand{\leq}{\leqslant}
\renewcommand{\ge}{\geqslant}
\renewcommand{\geq}{\geqslant}

\newcommand{\secpar}{\lambda}

%% bits

\newcommand{\bit}[1]{\ensuremath{\textcolor{bitcolor}{\texttt{\upshape #1}}}\xspace}
\newcommand{\bits}{\{\bit0,\bit1\}}

%% algorithms
\newcommand{\algorithm}[1]{\ensuremath{\textsf{\upshape#1}}\xspace}

\newcommand{\Scheme}{\Sigma}
\newcommand{\KeyGen}{\algorithm{KeyGen}}
\newcommand{\Enc}{\algorithm{Enc}}
\newcommand{\Encaps}{\algorithm{Encaps}}
\newcommand{\Sign}{\algorithm{Sign}}
\newcommand{\Dec}{\algorithm{Dec}}
\newcommand{\Decaps}{\algorithm{Decaps}}
\newcommand{\MAC}{\algorithm{MAC}}
\newcommand{\Verify}{\algorithm{Verify}}
\newcommand{\Share}{\algorithm{Share}}
\newcommand{\Reconstruct}{\algorithm{Reconstruct}}
\newcommand{\Filter}{\algorithm{Filter}}
\newcommand{\Birthday}{\algorithm{Birthday}}
\newcommand{\Start}{\algorithm{Start}}
\newcommand{\Respond}{\algorithm{Respond}}
\newcommand{\Finalize}{\algorithm{Finalize}}
\newcommand{\RSA}{\algorithm{RSA}}
\newcommand{\Commit}{\algorithm{Commit}}
\renewcommand{\Check}{\algorithm{Check}}
\newcommand{\Extract}{\algorithm{Extract}}
\newcommand{\SimTrans}{\algorithm{SimTranscript}}

\newcommand{\subname}[1]{\textsc{#1}}

%% misc crypto

\newcommand{\A}{\mathcal{A}}
\renewcommand{\L}{\mathcal{L}}
\newcommand{\Z}{\mathbb{Z}}
\newcommand{\Real}{\mathbb{R}}
\newcommand{\K}{\mathcal{K}}
\newcommand{\M}{\mathcal{M}}
\newcommand{\E}{\mathcal{E}}
\newcommand{\C}{\mathcal{C}}
\newcommand{\R}{\mathcal{R}}
\newcommand{\I}{\mathcal{I}}
\renewcommand{\S}{\mathcal{S}}
\newcommand{\T}{\mathcal{T}}
\newcommand{\X}{\mathcal{X}}
\newcommand{\Y}{\mathcal{Y}}
\newcommand{\randk}{\mathsf k}
\newcommand{\randm}{\mathsf m}
\newcommand{\randc}{\mathsf c}
\newcommand{\Adv}{\mathcal{A}}
\newcommand{\B}{\mathcal{B}}
\newcommand{\SSadv}{\textrm{SS}\textsf{adv}}
\newcommand{\SSkeyleak}{\textrm{SSKeyLeak}}
\newcommand{\SSkeyleakadv}{\textrm{SSKeyLeak}\textsf{adv}}


\newcommand{\msg}{M}
\newcommand{\key}{K}
\newcommand{\ctext}{C}

\newcommand{\attack}{\textsc{attack}}
\newcommand{\getrandom}{\twoheadleftarrow}
\newcommand{\indist}{  \mathrel{\vcenter{\offinterlineskip
  \hbox{$\sim$}\vskip-.35ex\hbox{$\sim$}\vskip-.35ex\hbox{$\sim$}}}}
\newcommand{\cenc}{\textsf{CaesarEnc}}
\newcommand{\cdec}{\textsf{CaesarDec}}
\newcommand{\rfrom}{\buildrel{\scriptstyle{\mathrm{R}}}\over{\leftarrow}}
\newcommand{\com}{\mathsf{Comp}}
\newcommand{\decom}{\mathsf{Decomp}}
\newcommand{\PRGadv}{\textrm{PRG}\textsf{adv}}
\newcommand{\countryguessadv}{\textsf{CountryGuessAdv}}
\newcommand{\country}{\mathsf{Country}}
\newcommand{\parity}{\mathsf{parity}}
\newcommand{\PRFadv}{\textrm{PRF}\textsf{adv}}
\DeclareMathOperator{\ord}{ord}
\newcommand{\MACadv}{\textrm{MAC}\textsf{adv}}
\newcommand{\CPAadv}{\textrm{CPA}\textsf{adv}}
\newcommand{\AltCPAadv}{\textrm{Alt-CPA}\textsf{adv}}
\newcommand{\BCadv}{\textrm{BC}\textsf{adv}}
\newcommand{\CRadv}{\textrm{CR}\textsf{adv}}
\newcommand{\checksum}{\textsf{Checksum}}

\newcommand{\samples}{\xleftarrow{R}}
\DeclareMathOperator{\Funs}{Funs}
\newcommand{\onehalf}{{\fontfamily{ppl}\selectfont\textonehalf}}
\newcommand{\onequarter}{{\fontfamily{ppl}\selectfont\textonequarter}}
\newcommand{\threequarters}{{\fontfamily{ppl}\selectfont\textthreequarters}}
\usepackage{amsthm,amsfonts,amsmath,xspace,tikz,varwidth,cancel,tikz}
\usetikzlibrary{arrows.meta, calc, chains, shapes.geometric, positioning, decorations.pathreplacing}

\begin{document}
\begin{center}
\Large{\textbf{CAS CS 538.  \ifsol Solutions to \fi Problem Set 6}}\\
\smallskip
\large{\textbf{Due \textcolor{red}{Thursday March 5}, 2026 11:59pm
}}\\
\large{\textbf{Om Khadka, U51801771}}
\end{center}

\noindent
\textbf{How to submit:}  See instructions on Problem Set 1. 


\begin{problem} (30 points total)
Let $F$ be a secure PRF with $n$-bit strings for keys, inputs, and outputs.

\begin{ppart} (5 points)
Let $F_1(k,x)$ computed as follows: compute $F(k, x)$, append the last bit of $x$, and output the result. Prove that $F_1$ is insecure.
\end{ppart}

\begin{solution}
The main idea behind this proof is the fact that the function, $F_1$ exposes a part of the input in its output. This can be exploited.

\subsection*{Making the Game}
Let $\A$ be the adversary that will play the game of distinguinish $F_1$ from just some rnd function.
\begin{itemize}
  \item $\A$ chooses any $x \in^R \{0, 1\}^n$
  \item Push $x$ into $\text{Oracle}$, producing $y \in \{0, 1\}^{n+1}$
  \item Now we check if $x[-1] = y[-1]$
  \begin{itemize}
    \item If YES, $\hat{b} = 1$
    \item Else, $\hat{b} = 0$
  \end{itemize}
  \item Return $\hat{b}$ to challenger.
\end{itemize}

\subsection*{Analysis}
Let $\PRFadv[\A, F_1] = |\Pr[W_0] - \Pr[W_1]|$ be the magnitude of difference of $\A$ outputting $1$ in either experiments $0, 1$.

\paragraph{$\exp 0$: Using $F_1$}
Assume that the challenger chooses $k \in^R \{0, 1\}^n$. In this experiment, the oracle is $F_1$. Then, by definition, the last bit of the output will ALWAYS be the last bit of $x$. Therefore $\hat{b} = 1$ always occurs.

\paragraph{$\exp 1$: Using rnd}
Now, the oracle is just a rnd func $R$. This means that every bit of $y$ is uniformly distributed and random to $x$'s bits. Thus, the chance of $\hat{b} = 1$ is $\frac{1}{2}$.

\subsection*{Adv}
Thus, the advantage for this game can be calculated as:
\begin{equation*}
  \PRFadv[\A, F_1] = |1 - \frac{1}{2}| = \frac{1}{2}
\end{equation*}
This is non-negligible, thus meaning that $F_1$ is insecure. $\qed$
\end{solution}

\begin{ppart} (5 points)
Let $F_2(k,x)$ computed as follows: compute $F(k, x)$, append the last bit of $k$, and output the result. Prove that $F_2$ is insecure.
\end{ppart}

\begin{solution}
This issue now is that the key is exposed in the output, which can be exploited.

\subsection*{Game}
$\A$ is playing to distinguish $F_2$ from a rnd function.
\begin{itemize}
  \item $\A$ chooses $x_1, x_2 \in^R \{0, 1\}^n$
  \item Push $x_1, x_2$ to $\text{Oracle}$ to get $y_1, y_2 \in \{0, 1\}^n$
  \item Now, check if $y_1[-1] == y_2[-1]$. This will set $\hat{b}$
  \begin{itemize}
    \item YES $\to 1$
    \item ELSE $\to 0$
  \end{itemize}
  \item Return $\hat{b}$ to $\A$
\end{itemize}

\subsection*{Analysis}
Let $\PRFadv[\A, F_2] = |\Pr[W_0] - \Pr[W_1]|$ be the magnitude of difference of $\A$ outputting $1$ in either experiments $0, 1$.

\paragraph{$\exp 0$: Using $F_2$}
Assume that the challenger chooses $k \in^R \{0, 1\}^n$. Denote $k_{\text{last}}$ to be the last bit of $k$. In this experiment, the oracle is $F_2$. Then, by definition, the last bit of outputs $y_i$ will ALWAYS equal to $k_{\text{last}}$. Therefore $\hat{b} = 1$ always occurs.

\paragraph{$\exp 1$: Using rnd}
Now, the oracle is just a rnd func $R$. This means that every bit of $y_i$ is uniformly distributed and random to $k$'s bits. Thus, the chance of $\hat{b} = 1$ is $\frac{1}{2}$.

\subsection*{Adv}
Thus, the advantage for this game can be calculated as:
\begin{equation*}
  \PRFadv[\A, F_2] = |1 - \frac{1}{2}| = \frac{1}{2}
\end{equation*}
This is non-negligible, thus meaning that $F_2$ is insecure. $\qed$
\end{solution}

\begin{ppart} (10 points)
Define $F_3$  for $2n$-bit keys, $2n$-bit inputs, and $2n$-bit outputs, as follows: 
\[
F_3((k_1, k_2), (x_1, x_2))=F(k_1, x_1)|| F(k_2, x_2)\,,
\]
where $||$ denotes concatenation of strings. Prove that $F_3$ is insecure.
\end{ppart}

\begin{solution}
  The main issue is that parts of the output are dependent on PARTS of the input. I can change just 1 part of the input, and not completely change the output.

\subsection*{Game}
$\A$ is playing to distinguish $F_3$ from a rnd function.
\begin{itemize}
  \item $\A$ chooses $x_1, x_2, x'_2 \in^R \{0, 1\}^n$, where $x_2 \neq x'_2$
  \item Push $(x_1, x_2)$ to $\text{Oracle}$ to get $y = y_1 || y_2$, with each $y_i \in \{0, 1\}^{n}$
  \item Push $(x_1, x'_2)$ to $\text{Oracle}$ to get $y' = y'_1 || y'_2$, with each $y'_i \in \{0, 1\}^{n}$
  \item Now, check if $y_1 == y'_1$. This will set $\hat{b}$
  \begin{itemize}
    \item YES $\to 1$
    \item ELSE $\to 0$
  \end{itemize}
  \item Return $\hat{b}$ to $\A$
\end{itemize}

\subsection*{Analysis}
Let $\PRFadv[\A, F_3] = |\Pr[W_0] - \Pr[W_1]|$ be the magnitude of difference of $\A$ outputting $1$ in either experiments $0, 1$.

\paragraph{$\exp 0$: Using $F_2$}
Assume that the challenger chooses $k_1, k_2 \in^R \{0, 1\}^n$. In this experiment, the oracle is $F_3$. Note how, in both outputs of $y$, we're using $x_1$. That then means that both $y_1, y'_1$ are dependent on $x_1$. Therefore, it should be the case that $y_1 = y'_1$, and $\hat{b} = 1$ always occurs.

\paragraph{$\exp 1$: Using rnd}
Now, the oracle is just a rnd func $R$. This then means that both outputs of both $y$'s bits are uniformly distributed and random to ANY of $x_i$'s bits. Thus, the chance of $\hat{b} = 1$ occuring is $\frac{1}{2^n}$.

\subsection*{Adv}
Thus, the advantage for this game can be calculated as:
\begin{equation*}
  \PRFadv[\A, F_2] = |1 - \frac{1}{2^n}| \approx 1
\end{equation*}
$\frac{1}{2^n}$ approaches negligiblity as $n$ increases. A negligible minus a non-negligible is non-neglibile (the result is a constant, in simpler terms). Thus, this is non-negligible, meaning that $F_3$ is insecure. $\qed$
\end{solution}


\begin{ppart} (10 points)
Define $F_4$  for $2n$-bit keys, $2n$-bit inputs, and $n$-bit outputs, as follows: 
\[
F_4((k_1, k_2), (x_1, x_2))=F(k_1, x_1)\oplus F(k_2, x_2)\,.
\]
Prove that $F_4$ is insecure.
\end{ppart}

\begin{solution}
The issue, like before, is the fact that sections of the output are dependent on sections of the input, allowing an adversary to modify inputs to find a pattern. It's a bit harder for this case tho, but XORs can be cancelled out.

\subsection*{Game}
$\A$ is playing to distinguish $F_4$ from a rnd function.
\begin{itemize}
  \item $\A$ chooses $x_1, x_2, x'_1, x'_2 \in^R \{0, 1\}^n$, where $x_1 \neq x'_1, x_2 \neq x'_2$
  \item Push the following 4 inputs into $\text{Oracle}$ to get the following $y_i \in \{0, 1\}^n$:
  \begin{itemize}
    \item $(x_1, x_2) \to \text{Oracle} \to y_1$
    \item $(x_1, x'_2) \to \text{Oracle} \to y_2$
    \item $(x'_1, x'_2) \to \text{Oracle} \to y_3$
    \item $(x'_1, x_2) \to \text{Oracle} \to y_4$
  \end{itemize}
  \item Now, check if $y_1 \oplus y_2 \oplus y_3 \oplus y_4 = \{0\}^n$. This will set $\hat{b}$
  \begin{itemize}
    \item YES $\to 1$
    \item ELSE $\to 0$
  \end{itemize}
  \item Return $\hat{b}$ to $\A$
\end{itemize}

\subsection*{Analysis}
Let $\PRFadv[\A, F_4] = |\Pr[W_0] - \Pr[W_1]|$ be the magnitude of difference of $\A$ outputting $1$ in either experiments $0, 1$.

\paragraph{$\exp 0$: Using $F_2$}
Assume that the challenger chooses $k_1, k_2 \in^R \{0, 1\}^n$. In this experiment, the oracle is $F_4$. Note the number of times each variables 
$x_1, x_2, x'_1, x'_2$ appears in the equation of $y_1 \oplus y_2 \oplus y_3 \oplus y_4$ ($y_i = (x_{\text{smth}}) \oplus (x_{\text{smth}})$). What you'll find is that, 
each of the following variable I have just listed out will come out EXACTLY TWICE.\\
\indent Also note how many times $k_1, k_2$ appears in the same equation. As we're XORing 4 terms, these 2 variables appear an EVEN NUMBER of times. 
By the properties of XOR, this then means that $y_1 \oplus y_2 \oplus y_3 \oplus y_4$ entirely cancels out to $0$. Therefore, $\hat{b} = 1$ always occurs.

\paragraph{$\exp 1$: Using rnd}
Now, the oracle is just a rnd func $R$. This means that each bit of each particular $y_i$ will be uniform and random to one another. However, I am not
necessarily saying that each $y_i$ needs to have the same exact bit. $y_1 \oplus y_2 \oplus y_3 = z$ is random, so you could just simlify to check if
$z \oplus y_4 = \{0\}^n$. As we're just checking if XORing a rnd value yields all $0$s, the chance of this happening is $\frac{1}{2^n}$.

\subsection*{Adv}
Thus, the advantage for this game can be calculated as:
\begin{equation*}
  \PRFadv[\A, F_4] = |1 - \frac{1}{2^n}| \approx 1
\end{equation*}
$\frac{1}{2^n}$ approaches negligiblity as $n$ increases. A negligible minus a non-negligible is non-neglibile (the result is a constant, in simpler terms). Thus, this is non-negligible, meaning that $F_4$ is insecure. $\qed$
\end{solution}
\end{problem}

\begin{problem} (20 points)
There are many PRF designs out there based on different assumptions. Suppose you aren't sure which ones are really secure and want to hedge your bets. In this problem, you will prove that you can take multiple different PRF designs and combine them  so that the result is secure as long as at least one them is secure. This type of risk reduction is called a PRF \emph{combiner}.  There are also combiners for other cryptographic primitives, like CPA-secure encryption, collision-resistant hashing, etc. In this problem, you will prove that XOR is a good PRF combiner.

Let $F_1$ and $F_2$ be functions with $n$-bit keys, inputs, and outputs. Let $F((k_1, k_2), x) = F_1(k_1, x)\oplus F_2(k_2, x)$. Prove that if either $F_1$ or $F_2$ is a secure PRF, then $F$ is a secure PRF.  (In fact, it suffices to prove that if $F_1$ is secure than $F$ is secure, because the proof for $F_2$ is the same. So just prove it for $F_1$ to get full credit.)

NB: your proof will, of course, proceed via a reduction. In your reduction, be sure to analyze two cases: when the $F_1$ adversary is in Experiment 0 of the PRF Attack Game (Game 4.2 in Section 4.4.1), you need to prove that the $F$ adversary is also seeing Experiment 0; when the $F_1$ adversary is in Experiment 1, you need to prove that the $F$ adversary is also seeing Experiment 1.
 \end{problem}

\begin{solution}
To solve this question, the following lemma is be essentially to proving my point:

\begin{lemma}[Random $\oplus$ Fixed]\label{lemma:rnd-xor-fixed}
  Let $R$ be a random function, and $C$ be \textbf{any} fixed function. Then, the following equation will also be random:
  \begin{equation*}
    H(x) = R(x) \oplus C(x)
  \end{equation*}
\end{lemma}
\begin{lem_explanation}
  Imagine $R$ being just a rnd bit ($0$ or $1$), and $C$ was $\{1\}$. Then, $R \oplus C$ outputting $1$ only occurs $50\%$ of the time (it is random).\\
  \indent Abstract this to an $n$-bit function for $R, C$ now. The probability of $R(x) \oplus C(x) = Z$, with $Z$ just being some
  desired $n$-bit output, will then also be random. This proves the lemma.
\end{lem_explanation}

\subsection*{Game}
\textit{Note:} I'll be proving the contrapositive: If $F$ is insecure, then $F_1$ is also insecure.\\
\indent Let $\B$ is the adversary playing against $F_1$, and $\A$ be the game playing against $F$ itself. I'll only really describe $\B$ in detail, $\A$ will be described more abstractly.
\begin{itemize}
  \item[\textbf{A}] $\A$ chooses some $x$
  \item $x$ gets pushed into $\B$
  \item[\textbf{B}] $x$ gets pushed into $\text{oracle}$, producing $y = F_1(k_1, x)$
  \item $\B$ chooses $k_2 = \{0, 1\}^n$
  \item $\B$ calculates $z = y \oplus F_2(k_2, x)$
  \item $\B$ returns $z$ to $\A$
  \item[\textbf{A, B}] $\A$ determines value of $\hat{b}$ from $z$
  \item $A, B$ returns $\hat{b}$
\end{itemize}

\subsection*{Analysis}
I won't explicity write out the $\Pr$ equation right now. HOWEVER note that experiment $0$ denotes the psuedo-rnd experiment, and $1$ denotes the truly random experiment for BOTH $\A, \B$.

\paragraph{Game 0 (For $\B$)}
In this game, the oracle $\B$ uses is $F_1$. Note how both terms of the value that $\B$ outputs to $\A$ are deterministic, $z = F_1(k_1, x) \oplus F_2(k_2, x)$. 
Further note how game we have just described is the exact same game as the problem asks of us. This means that the experiment $\A$ is playing then
is it's pseudo-rnd experiment as well. As $\B$'s output is dependent on $\A$, then:
\begin{equation*}
  \Pr[W^{\B}_0] = \Pr[W^{\A}_0]
\end{equation*}

\paragraph{Game 1 ($\B$)}
In this game, the orcle $\B$ now uses is just some rnd func. Now note the terms that $\B$ outputs to $\A$, $z = R(x) \oplus F_2(k_2, x)$. 
This is just \ref{lemma:rnd-xor-fixed}, which means that $z$ is a random value. That then means that the game $\A$ will play will then be its rnd experiment.
Thus:
\begin{equation*}
  \Pr[W^{\B}_1] = \Pr[W^{\A}_1]
\end{equation*}

\subsection*{Adv}
The advantage for game $\B$ can be calculated as the following:
\begin{equation*}
  \PRFadv[\B, F_1] = |\Pr[W^{\B}_0] - \Pr[W^{\B}_1]|
\end{equation*}

\noindent By our analysis, our probabilities are the exact same as $\A$'s probabilties:
\begin{equation*}
  |\Pr[W^{\B}_0] - \Pr[W^{\B}_1]| = |\Pr[W^{\A}_0] - \Pr[W^{\A}_1]| = \epsilon
\end{equation*}

\noindent As $\epsilon$ is assumed to be non-negligble (from our assumption), then this proves that $F_1$ is insecure, proving our contrapositive statement, and thus proving the OG statement. $\qed$
\end{solution}



\begin{problem} (20 points, at 10 each)

\begin{ppart}
  Let $F:\K\times \X \rightarrow \Y$ with $\K=\X=\Y=\{0,1\}^\lambda$ be a secure PRF. Define $F^1(k, x)$ is follows: if $k=0^\lambda$, output $0^\lambda$; else output $F(k, x)$. Prove that $F^1$ is secure. Hint: use the Difference Lemma (Theorem 4.7 in Section 4.4.3).
\end{ppart}

\begin{solution}
The hint does heavy lifting here. Basically this is a matter of showing how the probability of getting the `bad event' is negligble, which bounds to the overall probability of the PRF.

\begin{lemma}[Difference]\label{lemma:difference}
  Already learned it in class, so I'll just put the equation it states down here:
  \begin{equation*}
    |\Pr[W_0] - \Pr[W_1]| \leq \Pr[Z]
  \end{equation*}
\end{lemma}

\subsection*{Game}
  Define $\A$ as some abstract game that tries to distinguish $F^1$ from a truly rnd func.\\
  
  \noindent There are 3 total branches $\A$ could take, and it regards the oracle it uses:
  2 from the psuedo-rnd experiment, and 1 for the truly random:
  
  \begin{itemize}
    \item[Game 0] $F^1$ is run
    \item[Game 1] $F$ is run
    \item[Game 2] A rnd func is run  
  \end{itemize}

\subsection*{Analysis}
Denote $\PRFadv[\A, F] = |\Pr[W_0] - \Pr[W_1] - \Pr[W_2]|$ as the magnitude of probability difference of either game 0, 1, or 2 of occuring. By the triangle inequality,
this entire equation is bounded to the following expression:
\begin{equation*}
  \PRFadv[\A, F] \leq |\Pr[W_0] - \Pr[W_1]| + |\Pr[W_1] - \Pr[W_2]|
\end{equation*}
\noindent We can now analyze these 2 terms.

\paragraph{Game 0 $\to$ Game 1} Note how the literal difference in what $\A$ does between Game 0 and Game 1 is the exact same, EXCEPT when $k = 0^{\lambda}$.
Denote the event of this occuring as event $\Z$. Then, by \ref{lemma:difference}, the expression $|\Pr[W_0] - \Pr[W_1]|$ is bounded to $\Pr[\Z]$. 
Event $\Z$ only occurs when 1 particular element in $\K$ occurs (that is, $0^{\lambda}$). So then:
\begin{equation*}
  |\Pr[W_0] - \Pr[W_1]| \leq \Pr[\Z] \leq \frac{1}{2^{\lambda}}
\end{equation*}

\paragraph{Game 1 $\to$ Game 2} Game 1 is assumed to be secure. Game 2 uses a rnd oracle, so the output is uniform and rnd to any possible check $\A$ could do.
These are just 2 negligibles:
\begin{equation*}
  |\Pr[W_1] - \Pr[W_2]| <= \text{negligible}(\lambda)
\end{equation*}

\subsection*{Adv}
So then, the overall advantage of this question becomes bounded to the following:

\begin{equation*}
  \PRFadv[\A, F] \leq |\frac{1}{2^{\lambda}}| + |\text{negligible}(\lambda)|
\end{equation*}

\noindent Subtracting 2 negligibles still results in negligible. Thus, $F$ is secure. $\qed$

\end{solution}


 \begin{ppart} (10 points)
Use the previous part to prove the following.  Let $H:\K\times \X \rightarrow \Y$ with $\K=\X=\Y=\{0,1\}^\lambda$ be a secure PRF. Define $H^1(k', x')$ as follows: set $k=x'$ and  $x=k'$; output $H(k, x)$ (that is, simply swap the key and the output). Prove that $H^1$ is not necessarily a secure PRF.
 
 \end{ppart}
\begin{solution}
The issue here becomes two-fold. Firstly, we're using the input (smth the Adversary CAN control) as a conditional within the PRF.
Next is the weakness to the key, $k$. That weakness becomes visible here.

\subsection*{Game}
Let $\A$ be the adversary distinguishing $H'$ from a rnd func.

\begin{itemize}
  \item $\A$ chooses $x' \in \{0\}^{\lambda}$
  \item Push $(k', x') \to \text{Oracle} \to y \in \{0, 1\}^{\lambda}$
  \item Check now if $y == x'$. This is setting $\hat{b}$
  \begin{itemize}
    \item YES $\to$ 1
    \item NO $\to$ 0
  \end{itemize}
  \item Return $\hat{b}$ from $\A$
\end{itemize}

\subsection*{Analysis}
Let $\PRFadv[\A, H'] = |\Pr[W_1] - \Pr[W_2]|$ be the magnitude of difference between events $0, 1$ from outputting 1 from $\A$.

\paragraph{Game 0}  In this game, the oracle used is $H'(k', x') \to H(x', k')$. Note how the key here is now the input, $x'$, which the adversary set as $\{0\}^{\lambda}$.
By definition of $H(k, x)$, this then means that $y$ must also then be $0^{\lambda}$. Thus $\Pr[W_0] = 1$

\paragraph{Game 1}  In this game, a rnd func is used by oracle. As the outputs' bits are independent to the input, the probability is just $\frac{1}{2^{\lambda}}$.

\subsection*{Adv}
Thus the advantage is the following:

\begin{equation*}
  \PRFadv[\A, H'] = |1 - \frac{1}{2^{\lambda}}| \approx 1
\end{equation*}

\noindent A negligible minus a non-negligible is non-negligible (the difference is a const). Also, $\frac{1}{2^{\lambda}}$ approaches negligible as $\lambda$ increases.
Therefore, this proves that $H'$ is insecure. $\qed$
\end{solution}
 
\end{problem} 
 


\begin{problem} (30 points, at 15 each) 
Problem 5.3 from Section 5.8 of the Boneh + Shoup textbook. Specifically,

\begin{ppart}
Assume that $\E$ is CPA-secure and prove that $\E$ is Alt-CPA-secure.
\end{ppart}

\begin{ppart}
Assume that $\E$ is Alt-CPA-secure and prove that $\E$ is CPA-secure.
\end{ppart}


Hint 1: Do these parts one a time, to avoid getting lost. Part (b) requires a hybrid argument with many hybrid points.

Hint 2: The textbook presents the Alt-CPA definition as a single randomized experiment, with the bit $b$ chosen at random by the challenger, and the adversary is trying to guess b by outputting a matching $\hat b$. This definition style is called ``bit-guessing''; define 
\[
\AltCPAadv^*(\A, \E)=|\Pr[b=\hat b]-1/2|\,.
\]
This definition is equivalent to a two-experiment definition that we are more familiar with: one for $b=0$, and the other for $b=1$. In this view, there is no random choice of $b$, and we are not measuring whether $b=\hat b$; rather, we measure the probability difference between $p_0=\Pr[\hat b = 1]$ in Experiment $0$ and $p_1=\Pr[\hat b = 1]$ in Experiment 1; define 
\[
\AltCPAadv(\A, \E)=|p_1-p_0|\,.
\]
The two definitions are equivalent; this equivalence is shown in Section 2.2.5 in the textbook (in the context of semantic security, but it's the same for all other such definitions). It is easier to use $\CPAadv^*$ / $\AltCPAadv^*$ for part (a),  and $\CPAadv$ / $\AltCPAadv$ for part (b).
\end{problem}

\begin{solution}
\section*{A, $\CPAadv$}
The solution to part (a) of this question simply involves creating a reduction between $\AltCPAadv, \CPAadv$. In this case, $\A =$ the Alt-CPA fighter.
The goal is to make $\B$, the CPA fighter, perfectly simulate $\A$. \\

\noindent We try to solve the contrapositive: If $\E$ isn't Alt-CPA secure, then $\E$ isn't CPA secure.

\subsection*{Game}
Let $\A$ be the adversary fighting the Alt-CPA adversary. $\B$ is the adversary fighting the CPA adversary. $\A$ runs within $\B$.\\

\noindent Firstly, I'll describe $\A$:
\begin{itemize}
  \item IF $\A$ decides to run an enc query:
  \begin{itemize}
    \item $\A$ chooses $m \in \M$
    \item Pushes $m \to \text{Oracle}$
    \item Gets $c \in \Enc(k, m)$
    \item This whole loop can be run multiple times
  \end{itemize}
  \item IF $\A$ decides to run a test query:
  \begin{itemize}
    \item $\A$ chooses $m_1, m_2 \in \M$
    \item Pushes $(m_1, m_2) \to \text{Oracle}$
    \item Gets $c \in \Enc(k, m_b)$
    \item Can only be run once
  \end{itemize}
  \item Determines and returns $\hat{b}$
\end{itemize}

\noindent Note that, for $\A$, the oracle will ALWAYS be $\B$ here. Now, for the $\B$ game:
\begin{itemize}
  \item Gets $\text{query}$ from challenger.
  \item IF $\text{query} = m$ (a single var):
  \begin{itemize}
    \item Sends $(m, m) \to \text{oracle}$
    \item Gets $\Enc(k, m)$
    \item Sends $\Enc(k, m) \to$ challenger
  \end{itemize}
  \item IF $\text{query} = (m_1, m_2)$ (2-var tuple):
  \begin{itemize}
    \item Sends $(m_1, m_2) \to \text{oracle}$
    \item Gets $\Enc(k, m_b)$
    \item Sends $\Enc(k, m_b) \to$ challenger
  \end{itemize}
\end{itemize}

\noindent Note, $\text{oracle}$ for $\B$ is the CPA oracle, and $\text{challenger}$ is $\A$. Once $\A$ determines and returns $\hat{b}$, $\B$ returns the same value.

\subsection*{Analysis}
Let $\CPAadv[\B, \E] = |\Pr[W^{\B}_0] - \Pr[W^{\B}_1]|$ be the magnitude of difference of events $0, 1$ occuring within $\B$. \\

\noindent Also, let $\AltCPAadv[\A, \E] = |\Pr[W^{\A}_0] - \Pr[W^{\A}_1] = 2 \times \AltCPAadv^*[\A, \E]$ be the same for $\A$. As $\AltCPAadv^*$ is the bit-guessing advantage,
the distinguishing advantage is just half of that. For simplicity, let our inital assumption of that non-negligble advantage, $\epsilon$, be the advatange of
$\AltCPAadv$, not $\AltCPAadv^*$.

\paragraph{Game 0: CPA Oracle encrypts left} In the game, the $\text{oracle}$ $\B$ interacts with encrypts $m_0$. Consider the two cases of queries $\B$ could get:
\begin{itemize}
  \item[$m$] If $\A$ does an enc test, then $\B$ gets and sends $\Enc(m)$. Note how there's no difference between the bits from $\B$'s input, $(m, m)$. Thus, what $\A$
  gets from $\B$ is just the same as what $\B$ got. This is just game $\A$'s enc query.
  \item[$(m, m)$] In this case, now $\B$ gets $\Enc(m_0)$. Note how $\B$ sends this output as input to $\A$. This is just game $\A$'s test query.
\end{itemize}
\noindent Thus, the probability distribution of this event will be the same as what would also occur for game $\A$: $\Pr[W^{\B}_0] = \Pr[W^{\A}_0]$

\paragraph{Game 1: CPA Oracle encrypts right} In the game, the $\text{oracle}$ $\B$ interacts with encrypts $m_1$. Consider the two cases of queries $\B$ could get:
\begin{itemize}
  \item[$m$] If $\A$ does an enc test, then $\B$ gets and sends $\Enc(m)$. Like before, there's no difference between the outputs of $m, m$, they're the same.
  Thus, this is just game as $\A$'s enc query.
  \item[$(m, m)$] In this case, now $\B$ gets $\Enc(m_1)$. Note how $\B$ sends this output as input to $\A$. This is just game $\A$'s test query.
\end{itemize}
\noindent Same conclusion as before, just for Game 1 now: $\Pr[W^{\B}_1] = \Pr[W^{\A}_1]$

\subsection*{Adv}
Thus, the advantage for the CPA game is as follows:
\begin{equation*}
  \begin{aligned}
    \CPAadv &= |\Pr[W^{\B}_0] - \Pr[W^{\B}_1]|\\
    &= |\Pr[W^{\A}_0] - \Pr[W^{\A}_1]|\\
    &= \AltCPAadv[\A, \E]\\
    &= \epsilon
  \end{aligned}
\end{equation*}

\noindent Our assumpution assumed $\epsilon$ is non-negligble. Thus, $\CPAadv$ is non-negligible.\\

\noindent This proves the contrapositive, and thus the OG statement. $\qed$

\section*{B, $\AltCPAadv$}
This proof requires doing a bunch of hybrid arguments. The reason for this is because of what this question really is asking versus the previous one.
The previosu question was proving the equivalence of 2 different definitions. But here, I'm forced to now check if $\AltCPAadv$'s mechanism of single-query
is just as secure as the  standard $\CPAadv$ multi-query test.\\

\noindent I will prove the contrapositive: If $\E$ is \textbf{not} CPA-secure, then $\E$ is \textbf{not} Alt-CPA-secure.


\subsection*{Hybrid Games}
Let $\A$ be the adversary, that can make $Q$ test queries, and is assumed to have non-negligible advantage $\epsilon$ against $\E$. 
Now let $\B$ be the adversary, that can only make 1 query, that uses $\A$ to break Alt-CPA security.\\

\noindent Also, let $Q+1$ be defined as the hybrid games $H_0, H_1, \ldots, H_Q$. It is polynomially bounded:
\begin{itemize}
    \item \textbf{Hybrid $H_0$:} All $Q$ test queries use $b=0$ (encrypt $m_0^{(i)}$ for all $i$)
    \item \textbf{Hybrid $H_1$:} First query uses $b=1$, remaining $Q-1$ queries use $b=0$
    \item \textbf{Hybrid $H_2$:} First 2 queries use $b=1$, remaining $Q-2$ queries use $b=0$
    \item $\ldots$
    \item \textbf{Hybrid $H_Q$:} All $Q$ test queries use $b=1$ (encrypt $m_1^{(i)}$ for all $i$)
\end{itemize}

Note how $H_0$ corresponds to Experiment 0 of the CPA game, and how $H_Q$ corresponds to Experiment 1 of the CPA game.

\subsection*{Adversary}
$\B$ is now playing against the Alt-CPA game, and works as the CPA game's oracle within $\A$:
\begin{itemize}
    \item $\B$ chooses a random value $\omega \in \{1, \ldots, Q\}$
    \item When $\A$ makes its $i$-th test query $(m_0^{(i)}, m_1^{(i)})$:
    \begin{itemize}
        \item If $i < \omega$: $\B$ encrypts $m_1^{(i)}$ using its $\text{oracle}$ via an enc query
        \item If $i > \omega$: $\B$ will encrypt $m_0^{(i)}$ with $\text{oracle}$ via an enc query
        \item If $i = \omega$: $\B$ will finally send $(m_0^{(i)}, m_1^{(i)})$ to $\text{oracle}$ as a test query, get $c^*$, and send that to $\A$
    \end{itemize}
    \item $\B$ returns whatever $\A$ outputs as $\hat{b}$
\end{itemize}

\subsection*{Analysis}
Let $W_i$ be the event that $\A$ outputs 1 in the (hybrid) game $H_i$.

\paragraph{Case 1: $\B$'s oracle uses $b=0$}
When the Alt-CPA oracle encrypts $m_0^{(\omega)}$, note how $\B$'s game just becomes $H_{\omega-1}$ (as the first $\omega-1$ queries use $b=1$, while the rest use $b=0$).

\paragraph{Case 2: $\B$'s oracle uses $b=1$}
When the Alt-CPA oracle encrypts $m_1^{(\omega)}$, then this is just $H_{\omega}$ (the first $\omega$ queries use $b=1$, rest use $b=0$).

\subsection*{Adv}
The advantage of $\B$ against Alt-CPA thus is:
\begin{equation*}
\AltCPAadv[\B, \E] = \left|\Pr[\hat{b} = b] - \frac{1}{2}\right|
\end{equation*}

\noindent And since $\omega$ is independent and random from $\{1, \ldots, Q\}$:
\begin{equation*}
\AltCPAadv[\B, \E] \leq \frac{1}{Q} \sum_{i=1}^{Q} \left|\Pr[W_i] - \Pr[W_{i-1}]\right|
\end{equation*}

\noindent By the triangle inequality:
\begin{equation*}
\CPAadv[\A, \E] = \left|\Pr[W_Q] - \Pr[W_0]\right| \leq \sum_{i=1}^{Q} \left|\Pr[W_i] - \Pr[W_{i-1}]\right|
\end{equation*}

\noindent And therefore:
\begin{equation*}
\AltCPAadv[\B, \E] \geq \frac{1}{Q} \times \CPAadv[\A, \E]
\end{equation*}

\noindent Our initial assumption tells us that $\CPAadv$ will have some non-negligible value.
Since $Q$ is polynomially bounded (furthermore, it's $\frac{1}{Q}$, an inverse polynomial), that also then means that that product will still be non-negligible.\\

\noindent This then proves that $\AltCPAadv$ is non-negligible, and insecure. This proves the contrapositive, which proves our OG statment. $\qed$
\end{solution}
\end{document}